\documentclass[12pt]{article}
\usepackage[utf8]{inputenc}
\usepackage[T1]{fontenc}
\usepackage{hyperref}
\title{Urartian Language Resources for K-12 Curriculum Development}
\author{Compiled by Claude Sonnet 5.5 (claude-sonnet-5-5) for Andr\'as Kornai}
\date{}
\begin{document}
\maketitle

\section{Introduction}
Urartian (ISO 639-3 xur; Glottocode urar1245) is an extinct Hurro-Urartian language known from roughly two hundred cuneiform royal inscriptions of the Urartu kingdom (about the 9th to 6th centuries BCE). It has no living speakers and no community revitalization programme was found. Resources are scholarly.

\section{Primary Educational Resources}
No school textbooks or primers were found. Teaching material would have to be adapted from the scholarly resources below.

\section{Linguistic and Reference Materials}
\begin{itemize}
\item Wikipedia, ``Urartian language'': overview of the corpus, grammar and bibliography. \url{https://en.wikipedia.org/wiki/Urartian_language}
\item Palaeolexicon, ``The Urartian language (resources)'': a resource page with a bibliography of Urartian studies. \url{https://palaeolexicon.com/Urartian}
\item Leiden University repository item on Urartian (a scholarly text; I did not open it). \url{https://scholarlypublications.universiteitleiden.nl/access/item%253A3656156/view}
\item Journal of the American Oriental Society review of a recent concise grammar (co-written by Ilse Wegner), as surfaced by search. \url{https://lockwoodonlinejournals.com/index.php/jaos/article/download/1114/916}
\end{itemize}
The reference works named in search summaries, which I did not locate online as separate URLs, are: Salvini's \emph{Corpus dei testi urartei} (2018); the Electronic Corpus of Urartian Texts (eCUT, ed. Birgit Christiansen); Melikishvili's grammar (1963, German translation 1971); Ayvazyan's Urartian--Armenian lexicon (2011).

\section{Religious and Community Resources}
None found. The surviving texts are royal inscriptions, including dedications to the gods of Urartu, but no community organization for the language was found.

\section{Digital and Audio Resources}
Audio is not possible for an extinct language known only from cuneiform. The Palaeolexicon page and eCUT are the digital resources found. I did not verify eCUT's URL or availability.

\section{Curriculum Development Context}
The corpus is small. At K-12 level, a realistic use is a cultural and historical unit (Urartu, cuneiform, Armenian highland history) with a few attested words and royal-inscription excerpts, in translation. Using Urartian as the language of instruction is not feasible.

\end{document}